\section{补充证明与代码说明}

\subsection{问题一条件性互补化的交换证明}
\label{app:p1-relaxation-proof}
\begingroup\small
设$(g_t,c_t,d_t,w_t,E_t)$是命题1所述LP的一个最优解，且某时段$t$满足
$c_t>0,d_t>0$。记$\alpha=\eta_c\eta_d$，并定义
\[
 \delta_t=\min\left\{c_t,\frac{d_t}{\alpha}\right\},\qquad
 r_t=(1-\alpha)\delta_t.
\]
由于$\delta_t>0$，令
\[
 c'_t=c_t-\delta_t,\qquad d'_t=d_t-\alpha\delta_t
\]
后，$c'_t,d'_t\ge0$且至少一个等于0；两个变量只减不增，故原充、放电功率上限仍成立。
该变换对储能增量的影响为
\[
 \eta_c(c'_t-c_t)\Delta t-
 \frac{(d'_t-d_t)\Delta t}{\eta_d}
 =-\eta_c\delta_t\Delta t+
 \frac{\alpha\delta_t\Delta t}{\eta_d}=0.
\]
因此$E_{t+1}$不变，继而全部后续储能状态、SOC上下界及终端条件均保持不变。

在母线功率平衡中，上述动作产生$r_t$的供给余量。取
\[
 a_t=\min\{g_t,r_t\},\qquad b_t=r_t-a_t.
\]
由条件\eqref{eq:p1-exchange-condition}可得$0\le b_t\le P_t-w_t$。再令
\[
 g'_t=g_t-a_t,\qquad w'_t=w_t+b_t,
\]
则$g'_t\ge0$且$0\le w'_t\le P_t$。代入式\eqref{eq:p1-balance}，有
\[
 (g'_t+P_t+d'_t)-(L_t+c'_t+w'_t)
 =r_t-a_t-b_t=0,
\]
故功率平衡仍成立。其他时段变量保持不变，因而得到一个消除了时段$t$同时充放电的LP可行解。
目标函数的变化量为
\[
 \Delta z=-p_ta_t\Delta t\le0.
\]
对所有同时充放电时段重复该交换，得到满足$c_td_t=0$的LP可行解。由于原解已经最优，新解的
目标值不可能低于LP最优值，故交换后的目标值与原值相同。为该解补充二元变量：充电时取
$u_t=1$，放电时取$u_t=0$，两者均为0时任取，即得到同目标值的MILP可行解。因此
$z_{\mathrm{MILP}}\le z_{\mathrm{LP}}$；又因LP是MILP的松弛，
$z_{\mathrm{LP}}\le z_{\mathrm{MILP}}$，从而$z_{\mathrm{LP}}=z_{\mathrm{MILP}}$。

需要强调的是，条件\eqref{eq:p1-exchange-condition}保证消去有损循环后产生的供给余量能够通过
减少购电或增加弃光吸收。若不能验证这一条件，上述局部交换并不自动构成一般性的精确松弛证明；
附件1实例的零间隙结论仍以表\ref{tab:p1-lp-relaxation}的独立数值交叉验证为准。
\endgroup

\section{支撑材料与程序说明}

\subsection{支撑材料文件列表}
\begingroup\tiny
\setlength{\tabcolsep}{2pt}
\setlength{\LTleft}{0pt}\setlength{\LTright}{0pt}
\begin{longtable}{@{}C{0.04\textwidth}p{0.20\textwidth}C{0.065\textwidth}C{0.085\textwidth}p{0.175\textwidth}p{0.12\textwidth}p{0.18\textwidth}@{}}
\toprule
编号 & 文件路径 & 类型 & 对应问题 & 功能 & 输入 & 输出 \\
\midrule
\endfirsthead
\toprule
编号 & 文件路径 & 类型 & 对应问题 & 功能 & 输入 & 输出 \\
\midrule
\endhead
1 & {\raggedright data/\allowbreak{}raw/\allowbreak{}C题/\allowbreak{}附件/\allowbreak{}附件1.xlsx\par} & Excel & 问题一、二、三 & 固定电价、典型负荷与光伏 & 题目附件 & 确定性与固定价调度输入 \\
2 & {\raggedright data/\allowbreak{}raw/\allowbreak{}C题/\allowbreak{}附件/\allowbreak{}附件2.xlsx\par} & Excel & 问题二、三、四 & 全年实际负荷与光伏 & 题目附件 & 实际回放输入 \\
3 & {\raggedright data/\allowbreak{}raw/\allowbreak{}C题/\allowbreak{}附件/\allowbreak{}附件3.xlsx\par} & Excel & 问题三、四 & 多发布时间光伏预报 & 题目附件 & 滚动预测输入 \\
4 & {\raggedright data/\allowbreak{}raw/\allowbreak{}C题/\allowbreak{}附件/\allowbreak{}附件4.xlsx\par} & Excel & 问题四 & 10分钟动态电价 & 题目附件 & 动态价格向量 \\
5 & {\raggedright src/\allowbreak{}common/\allowbreak{}\par} & Python & 全部 & 数据校验、预处理、指标与绘图公共模块 & 附件1--4 & 规范化数据与诊断 \\
6 & {\raggedright src/\allowbreak{}forecasting/\allowbreak{}\par} & Python & 问题二 & 因果预测、SSA、场景生成与冻结评价 & 规范化源荷数据 & 预测与联合场景 \\
7 & {\raggedright src/\allowbreak{}problem1/\allowbreak{}\par} & Python & 问题一 & 确定性MILP、LP松弛与灵敏度 & 附件1 & result1及问题一图表 \\
8 & {\raggedright src/\allowbreak{}problem2/\allowbreak{}\par} & Python & 问题二 & 两阶段随机MILP、物理结算与CVaR & 附件1、2及预测场景 & result2、审计与风险表 \\
9 & {\raggedright src/\allowbreak{}problem3/\allowbreak{}\par} & Python & 问题三 & 预测融合、滚动重优化、结算与VOI & 附件1--3 & result3、S0--S3及VOI \\
10 & {\raggedright src/\allowbreak{}problem4/\allowbreak{}\par} & Python & 问题四 & 动态价格接口、年度调度与鲁棒性 & 附件2--4 & result4-2、result4-3及鲁棒性结果 \\
11 & {\raggedright src/\allowbreak{}analysis/\allowbreak{}information\_\allowbreak{}bound.py\par} & Python & 综合分析 & 离线全信息下界与信息差距 & 全年实际源荷 & 下界费用与审计 \\
12 & {\raggedright src/\allowbreak{}paper\_\allowbreak{}figures/\allowbreak{}\par} & Python & 全部 & 论文正式图生成 & 各问题结果表 & paper/figures中的正式图 \\
13 & {\raggedright results/\allowbreak{}problem1/\allowbreak{}result1.xlsx\par} & Excel & 问题一 & 题目格式的确定性调度结果 & 问题一模型 & 提交表1、表2 \\
14 & {\raggedright results/\allowbreak{}problem2/\allowbreak{}result2.xlsx\par} & Excel & 问题二 & 334日计划、储能与紧急购电结果 & 问题二模型 & 提交表3 \\
15 & {\raggedright results/\allowbreak{}problem3/\allowbreak{}result3.xlsx\par} & Excel & 问题三 & S3正式滚动购电与调整结果 & 问题三模型 & 问题三提交结果 \\
16 & {\raggedright results/\allowbreak{}problem4/\allowbreak{}result4-2.xlsx\par} & Excel & 问题四 & 动态价格下问题二重算结果 & Q4-2模型 & 问题4-2提交结果 \\
17 & {\raggedright results/\allowbreak{}problem4/\allowbreak{}result4-3.xlsx\par} & Excel & 问题四 & 动态价格下问题三重算结果 & Q4-3模型 & 问题4-3提交结果 \\
18 & {\raggedright results/\allowbreak{}problem2/\allowbreak{}Q2\_\allowbreak{}PHYSICAL\_\allowbreak{}SETTLEMENT\_\allowbreak{}AUDIT.md\par} & Markdown & 问题二 & 实际执行与成本口径审计 & Q2逐时结果 & 物理与费用核验 \\
19 & {\raggedright results/\allowbreak{}problem3/\allowbreak{}Q3\_\allowbreak{}ROLLING\_\allowbreak{}AUDIT.md\par} & Markdown & 问题三 & 因果性、冻结、结算及VOI审计 & Q3年度结果 & 问题三核验 \\
20 & {\raggedright results/\allowbreak{}problem4/\allowbreak{}Q4\_\allowbreak{}GIVEN\_\allowbreak{}PRICE\_\allowbreak{}INDEPENDENT\_\allowbreak{}AUDIT.json\par} & JSON & 问题四 & 动态价格年度独立审计 & Q4年度结果 & 价格、物理与结算核验 \\
21 & {\raggedright requirements-lock.txt\par} & 文本 & 全部 & Python依赖精确版本 & 环境冻结 & 可复现安装清单 \\
22 & {\raggedright paper/\allowbreak{}AI工具使用详情.tex\par} & LaTeX & 支撑材料 & AI工具使用详情源文件 & 真实使用记录 & AI工具使用详情.pdf \\
\bottomrule
\end{longtable}
\endgroup

\section{核心源程序}

\lstset{
  basicstyle=\ttfamily\scriptsize,
  breaklines=true,
  columns=fullflexible,
  keepspaces=true,
  numbers=left,
  numberstyle=\tiny,
  numbersep=5pt,
  showstringspaces=false
}

\subsection{数据预处理与场景构造}

\lstinputlisting[caption={\texttt{src/common/preprocessing.py}}]{../src/common/preprocessing.py}

\lstinputlisting[caption={\texttt{src/forecasting/final\_models.py}}]{../src/forecasting/final_models.py}

\lstinputlisting[caption={\texttt{src/forecasting/scenarios.py}}]{../src/forecasting/scenarios.py}

\subsection{问题一与问题二优化模型}

\lstinputlisting[caption={\texttt{src/problem1/model.py}}]{../src/problem1/model.py}

\lstinputlisting[caption={\texttt{src/problem2/model.py}}]{../src/problem2/model.py}

\lstinputlisting[caption={\texttt{src/problem2/forecast\_interface.py}}]{../src/problem2/forecast_interface.py}

\subsection{结果文件的时间标签与导出}

\lstinputlisting[firstline=30,lastline=69,caption={\texttt{src/problem1/run.py}中的时间标签与结果导出}]{../src/problem1/run.py}

\lstinputlisting[firstline=82,lastline=88,caption={\texttt{src/problem2/run.py}中的10分钟区间标签}]{../src/problem2/run.py}

\lstinputlisting[firstline=330,lastline=390,caption={\texttt{src/problem2/run.py}中的结果工作簿导出}]{../src/problem2/run.py}

\lstinputlisting[firstline=154,lastline=258,caption={\texttt{src/paper\_figures/specified\_date\_tables.py}中的题面格式表格生成}]{../src/paper_figures/specified_date_tables.py}

\subsection{问题三滚动优化}

\lstinputlisting[caption={\texttt{src/problem3/forecast\_fusion.py}}]{../src/problem3/forecast_fusion.py}

\lstinputlisting[caption={\texttt{src/problem3/rolling\_model.py}}]{../src/problem3/rolling_model.py}

\lstinputlisting[caption={\texttt{src/problem3/rolling\_dispatch.py}}]{../src/problem3/rolling_dispatch.py}

\subsection{问题四动态电价接口}

\lstinputlisting[caption={\texttt{src/problem4/price\_interface.py}}]{../src/problem4/price_interface.py}

\lstinputlisting[caption={\texttt{src/problem4/dispatch\_adapter.py}}]{../src/problem4/dispatch_adapter.py}

